\chapter{Mobile Execution Model}
\label{ch:mobile-execution}

The gert v2 design supports offline, on-device runbook execution on mobile platforms through an embedded SDK and platform-specific kit bundles. This chapter specifies the mobile execution architecture, the kit bundle format, per-platform tool dispatch, dependency resolution, and run synchronisation.

% ------------------------------------------------------------------ %
\section{Motivation and Design Goals}
% ------------------------------------------------------------------ %

The mobile execution model is designed to satisfy three requirements:

\begin{enumerate}
  \item \textbf{Offline-first execution:} Runbooks execute locally on-device without requiring network connectivity. The gert engine runs entirely within the mobile application process. Network access is only required when a specific runbook step explicitly requires it (e.g., calling an API), not for execution control flow.
  \item \textbf{Server-optional sync:} Evidence and trace data are collected locally. Synchronisation to a central gert server is opt-in and asynchronous. Organisations that require central audit trails can deploy a gert server and configure the SDK to sync completed runs; others can operate purely on-device.
  \item \textbf{Format compatibility:} The runbook YAML format, trace JSONL structure, evidence model, and state snapshot contract remain unchanged. Mobile execution is an alternative executor, not an alternative schema. Desktop-authored runbooks compile to mobile bundles without manual translation.
\end{enumerate}

The core principle is simple: \emph{gert runbooks run locally everywhere}. The format is 100\% compatible across CLI, server, and mobile platforms. Only the tool runtime adapter changes to dispatch platform-specific native implementations.

% ------------------------------------------------------------------ %
\section{Architecture Overview}
% ------------------------------------------------------------------ %

\subsection{Embedded Engine}

Mobile applications embed a gert engine instance via the GertSDK library. The SDK is distributed as:

\begin{itemize}
  \item \textbf{iOS:} Swift Package Manager package \texttt{GertSDK} exposing Swift APIs.
  \item \textbf{Android:} Maven/Gradle library \texttt{com.gert:gert-sdk} exposing Kotlin/Java APIs.
\end{itemize}

Both SDKs wrap a common gert core (written in Go and compiled to native libraries via \texttt{gomobile}). The engine instantiated within the mobile app process behaves identically to the CLI executor: it parses runbook YAML, evaluates branches, enforces governance, captures evidence, and writes trace JSONL.

\subsection{Kit Bundles}

Runbooks are packaged as \textbf{kit bundles} for mobile deployment. A kit bundle is a directory archive containing:

\begin{itemize}
  \item A manifest declaring the kit version, target platform, and required dependencies.
  \item Pre-lowered runbook YAML files (using only core gert primitives).
  \item Shared tool definitions scoped to the platform (iOS/Android).
  \item Catalog entries enabling discovery of runbooks within the bundle.
  \item Optional reference data (schemas, lookup tables) shared across runbooks.
\end{itemize}

Bundles are distributed via application bundle embedding, over-the-air update channels, or direct download from a registry. The SDK resolves bundle dependencies (platform kits, domain kits) at load time.

\subsection{Platform Kits}

A \textbf{platform kit} is a specialised kit bundle that provides mobile-specific tool implementations. The canonical platform kit, \texttt{gert-mobile-platform}, declares tools for:

\begin{itemize}
  \item \texttt{camera.capture} — capture photo evidence (iOS/Android native camera APIs)
  \item \texttt{location.get} — retrieve device GPS coordinates
  \item \texttt{nfc.read} — read NFC tags (equipment identifier scanning)
  \item \texttt{biometrics.verify} — authenticate operator via Touch ID / Face ID / fingerprint
  \item \texttt{bluetooth.scan} — discover nearby BLE devices
  \item \texttt{notifications.send} — send local push notifications
\end{itemize}

Third-party organisations can publish their own platform kits (e.g., \texttt{acme-enterprise-platform}) declaring custom tools specific to their environment. The compiler and SDK both reference the registry to resolve platform kit dependencies during bundle creation and loading.

% ------------------------------------------------------------------ %
\section{Kit Bundle Format}
% ------------------------------------------------------------------ %

A kit bundle is a directory with the following canonical structure:

\begin{verbatim}
<kit-name>.kit/
  manifest.json       # version, target platform, requires, checksum, signature
  catalog/            # one catalog-entry.yaml per runbook
    pool-maintenance.yaml
    lawn-mow.yaml
    irrigation-check.yaml
  runbooks/           # lowered runbook YAMLs (core primitives only)
    pool-maintenance.runbook.yaml
    lawn-mow.runbook.yaml
    irrigation-check.runbook.yaml
  tools/              # shared tools declared once for all runbooks
    camera.capture.tool.yaml
    location.get.tool.yaml
  assets/             # shared schemas, reference data
    equipment-schema.json
    region-lookup.csv
\end{verbatim}

\subsection{Manifest Schema}

The \texttt{manifest.json} file is the only new schema artefact introduced by the mobile execution model. All other files reuse existing gert v2 schemas.

\begin{minted}{json}
{
  "$schema": "https://schemas.gert.dev/kit-manifest/v1.json",
  "name": "home-domain",
  "version": "1.2.0",
  "target": "mobile",               // "mobile" | "ios" | "android"
  "requires": [
    {
      "kit": "gert-mobile-platform",
      "version": ">=2.0.0 <3.0.0"
    }
  ],
  "checksum": {
    "sha256": "a3f9c1...",            // SHA256 of all files excluding manifest
    "algorithm": "sha256"
  },
  "signature": {
    "keyId": "acme-signing-key-2026",
    "value": "base64-encoded-signature"
  },
  "metadata": {
    "description": "Home maintenance domain kit",
    "author": "Acme Corp",
    "published": "2026-04-18T10:00:00Z"
  }
}
\end{minted}

\paragraph{Target platform values:}

\begin{itemize}
  \item \texttt{"mobile"} — compatible with both iOS and Android (all tool impls present for both platforms)
  \item \texttt{"ios"} — iOS-only bundle
  \item \texttt{"android"} — Android-only bundle
\end{itemize}

The \texttt{requires} field declares dependency on other kits (typically platform kits). Version constraints use npm-style semver ranges.

\subsection{Catalog Entries}

Each runbook in the bundle has a corresponding catalog entry under \texttt{catalog/}. The catalog entry provides metadata for runbook discovery and search within the mobile application UI.

\begin{minted}{yaml}
# catalog/pool-maintenance.yaml
$schema: "https://schemas.gert.dev/catalog-entry/v1.yaml"
id: pool-maintenance
name: "Weekly Pool Maintenance"
description: "Check chemical levels, clean filters, inspect equipment"
tags: [pool, maintenance, weekly]
requires-capabilities: [capability/camera]
estimated-duration: "20m"
runbook: "runbooks/pool-maintenance.runbook.yaml"
\end{minted}

The \texttt{runbook} field is a relative path within the bundle. The SDK resolves it at load time.

% ------------------------------------------------------------------ %
\section{Platform Kit Concept and Registry}
% ------------------------------------------------------------------ %

\subsection{Hierarchy Model}

Kits are organised in a two-tier hierarchy:

\begin{verbatim}
Platform Kit (gert-mobile-platform — shared mobile capability tools)
  └── Domain Kit (e.g., home-domain)
        ├── pool-maintenance
        ├── lawn-mow
        └── irrigation-check
\end{verbatim}

The platform kit provides low-level tools that interact with device capabilities (camera, GPS, NFC). Domain kits depend on the platform kit and provide runbooks specific to a business domain (home maintenance, field service, asset inspection).

\subsection{Dependency Resolution}

Dependency resolution is \textbf{eager}. When the application calls \texttt{GertSDK.loadKit(bundlePath)}, the SDK:

\begin{enumerate}
  \item Reads \texttt{manifest.json} from the bundle.
  \item Resolves all \texttt{requires:} entries immediately.
  \item Downloads missing platform kits from the configured registry (if network is available).
  \item Validates all tool implementations for the target platform (fail early if any required tool lacks an \texttt{impl} block for the current platform).
  \item Registers all tools, runbooks, and catalog entries into the SDK's in-memory registry.
\end{enumerate}

If a required platform kit is missing and the device is offline, \texttt{loadKit} fails immediately with a dependency error. The application layer is responsible for bundling required platform kits in the application binary or ensuring they are pre-downloaded.

\subsection{Platform Kit Registry}

The canonical platform kit registry is hosted at \texttt{https://kits.gert.dev/}. Third-party organisations can self-host registries and configure the SDK to reference them:

\begin{minted}{json}
{
  "registries": [
    "https://kits.gert.dev/",
    "https://kits.acme.example/internal"
  ]
}
\end{minted}

Registry resolution follows the order declared. The SDK queries each registry in turn until the required kit is found.

% ------------------------------------------------------------------ %
\section{Per-Platform Tool Dispatch}
% ------------------------------------------------------------------ %

\subsection{Impl Blocks in Tool Definitions}

Mobile tool definitions extend the v2 \texttt{.tool.yaml} schema (Chapter~\ref{ch:tools}) with an optional \texttt{impl:} block declaring platform-specific handlers.

\begin{minted}{yaml}
# tools/camera.capture.tool.yaml
apiVersion: tool/v2

meta:
  name: camera.capture
  version: "1.0.0"
  description: "Capture photo evidence using device camera"

governance:
  requires-capabilities: [capability/camera]

actions:
  capture:
    description: "Capture a photo"
    args:
      evidence_name:
        type: string
        required: true
        description: "Name for the evidence item"
    output:
      format: json

impl:
  ios:
    transport: native-sdk
    handler: GertSDK.Camera.capture
  android:
    transport: native-sdk
    handler: com.gert.tools.camera.CaptureHandler
\end{minted}

The \texttt{impl} block maps platform identifiers (\texttt{ios}, \texttt{android}) to handler implementations. The \texttt{transport: native-sdk} value instructs the runtime to dispatch the tool call to a registered native handler class instead of spawning a process.

\subsection{Handler Registration (iOS)}

iOS handlers are Swift classes conforming to the \texttt{GertToolHandler} protocol:

\begin{minted}{go}
public protocol GertToolHandler {
    func invoke(action: String, args: [String: Any], 
                context: GertInvocationContext) async throws 
                -> GertToolResult
}

public class CameraCapture: GertToolHandler {
    public func invoke(action: String, args: [String: Any], 
                      context: GertInvocationContext) async throws 
                      -> GertToolResult {
        // Use AVFoundation to capture photo
        // Return evidence attachment reference
    }
}
\end{minted}

Handlers are registered at SDK initialisation:

\begin{minted}{go}
GertSDK.registerToolHandler(
    name: "camera.capture",
    handler: CameraCapture()
)
\end{minted}

\subsection{Handler Registration (Android)}

Android handlers implement the \texttt{GertToolHandler} interface:

\begin{minted}{java}
public interface GertToolHandler {
    GertToolResult invoke(String action, 
                         Map<String, Object> args,
                         GertInvocationContext context) 
                         throws Exception;
}

public class CaptureHandler implements GertToolHandler {
    @Override
    public GertToolResult invoke(String action, 
                                Map<String, Object> args,
                                GertInvocationContext context) {
        // Use CameraX API to capture photo
        // Return evidence attachment reference
    }
}
\end{minted}

Handlers are registered at SDK initialisation:

\begin{minted}{java}
GertSDK.registerToolHandler(
    "camera.capture", 
    new CaptureHandler()
);
\end{minted}

\subsection{Capability Gating at Load Time}

If a tool declares an \texttt{impl} block but lacks an entry for the target platform, the SDK emits a \texttt{capability/unavailable} error at \textbf{load time} (during \texttt{loadKit}), not at runtime.

\begin{verbatim}
ERROR: tool 'camera.capture' requires impl for platform 'ios' but none 
       is declared in tools/camera.capture.tool.yaml
\end{verbatim}

This fail-early policy ensures that runbook execution never begins if required tools cannot be dispatched. It is the mobile equivalent of the capability gate described in Section~\ref{sec:tool-capability-gate}.

% ------------------------------------------------------------------ %
\section{Bundle Hierarchy}
% ------------------------------------------------------------------ %

A typical mobile deployment consists of:

\begin{enumerate}
  \item One or more \textbf{platform kits} providing low-level tool implementations (camera, GPS, NFC). These are typically distributed by the gert project or the organisation's platform team.
  \item One or more \textbf{domain kits} providing business-specific runbooks (field service, asset inspection, maintenance checklists). These depend on platform kits and are authored by domain teams.
  \item The mobile application binary, which embeds or downloads these kits at runtime.
\end{enumerate}

The SDK allows multiple domain kits to coexist. The application layer presents a unified runbook catalog combining all loaded kits.

% ------------------------------------------------------------------ %
\section{Target Platform Compilation}
% ------------------------------------------------------------------ %

\subsection{Compiler Target Flag}

The gert compiler supports a \texttt{--target} flag to produce platform-specific bundles:

\begin{verbatim}
gert compile --target ios       --output home-domain.kit/
gert compile --target android   --output home-domain.kit/
gert compile --target mobile    --output home-domain.kit/
\end{verbatim}

The compiler performs the following transformations:

\begin{enumerate}
  \item \textbf{Lowering:} High-level runbook constructs (loops, conditionals, invoke) are lowered to core gert primitives that the mobile runtime supports.
  \item \textbf{Tool validation:} Every tool referenced by a runbook step is checked for an \texttt{impl} block matching the target platform. Missing impls produce a compile error.
  \item \textbf{Dependency bundling:} Required platform kits are resolved from the registry and their tool definitions are included in the bundle's \texttt{tools/} directory.
  \item \textbf{Manifest generation:} The \texttt{manifest.json} file is generated with correct \texttt{target}, \texttt{version}, \texttt{requires}, and \texttt{checksum} fields.
\end{enumerate}

\subsection{Platform-Specific Exclusions}

If a runbook step references a tool that lacks an \texttt{impl} for the target platform, the compiler behaviour depends on the step's \texttt{if:} condition:

\begin{itemize}
  \item If the step has no \texttt{if:} condition (unconditional step), the compiler emits an error and compilation fails.
  \item If the step is guarded by \texttt{if: platform == "ios"}, the compiler excludes the step when compiling for Android (and vice versa).
\end{itemize}

This allows runbook authors to write cross-platform runbooks with platform-specific branches.

% ------------------------------------------------------------------ %
\section{Run Sync and Ingest API}
% ------------------------------------------------------------------ %

\subsection{Offline Evidence Collection}

During runbook execution, the mobile SDK writes all trace events, snapshots, and attachments to local storage using the same directory structure as the CLI executor (see Section~\ref{sec:trace-file-format}):

\begin{verbatim}
<app-data-dir>/.runbook/runs/<run-id>/
  trace.jsonl
  snapshots/step-0000.json
  attachments/<sha256>.bin
  run.yaml
\end{verbatim}

All evidence is stored locally on-device. No network communication is required during execution.

\subsection{Run Ingest API}

When network connectivity is available and the application is configured to sync runs to a central server, the SDK uploads completed runs via the \textbf{run ingest API}.

\subsubsection{Endpoint: \texttt{POST /api/v1/runs/ingest}}

Accepts a stream of JSONL trace events (identical format to \texttt{trace.jsonl}). The server reconstructs the run and stores it in its run database.

\paragraph{Request format:}

\begin{verbatim}
POST /api/v1/runs/ingest
Content-Type: application/x-ndjson
Authorization: Bearer <token>

{"event_id":"...","sequence":0,"kind":"run/started",...}
{"event_id":"...","sequence":1,"kind":"step/started",...}
{"event_id":"...","sequence":2,"kind":"step/completed",...}
...
\end{verbatim}

\paragraph{Response:}

\begin{minted}{json}
{
  "run_id": "20260418T143022-a3f9c1",
  "status": "ingested",
  "events_received": 47,
  "attachments_pending": ["a3f9c1...", "b2e4d5..."]
}
\end{minted}

The \texttt{attachments\_pending} field lists SHA256 digests of attachments referenced in the trace but not yet uploaded.

\subsubsection{Endpoint: \texttt{POST /api/v1/runs/\{run-id\}/attachments/\{sha256\}}}

Uploads a single attachment file.

\paragraph{Request format:}

\begin{verbatim}
POST /api/v1/runs/20260418T143022-a3f9c1/attachments/a3f9c1...
Content-Type: application/octet-stream
Authorization: Bearer <token>

<binary file content>
\end{verbatim}

\paragraph{Response:}

\begin{minted}{json}
{
  "sha256": "a3f9c1...",
  "size": 204800,
  "status": "stored"
}
\end{minted}

\subsection{Sync Strategy}

The SDK provides two sync strategies:

\begin{itemize}
  \item \textbf{Immediate sync:} Upload trace and attachments immediately after \texttt{run/completed} event is written (requires active network).
  \item \textbf{Deferred sync:} Queue completed runs locally and upload when network becomes available (offline-first mode).
\end{itemize}

The application layer controls sync policy via SDK configuration:

\begin{minted}{json}
{
  "sync": {
    "enabled": true,
    "strategy": "deferred",
    "server_url": "https://gert.acme.example",
    "auth_token_provider": "app-delegate"
  }
}
\end{minted}

% ------------------------------------------------------------------ %
\section{SDK Responsibilities}
% ------------------------------------------------------------------ %

\subsection{iOS SDK API Surface}

\begin{minted}{go}
public class GertSDK {
    // Initialise SDK with configuration
    public static func configure(config: GertConfig)
    
    // Load a kit bundle from filesystem or network
    public static func loadKit(bundlePath: String) async throws
    
    // List all runbooks from loaded kits
    public static func listRunbooks() -> [RunbookCatalogEntry]
    
    // Execute a runbook
    public static func executeRunbook(
        id: String, 
        inputs: [String: Any],
        delegate: GertExecutionDelegate
    ) async throws -> GertRunResult
    
    // Resume a suspended run
    public static func resumeRun(runId: String) async throws 
                                 -> GertRunResult
    
    // Register a native tool handler
    public static func registerToolHandler(name: String, 
                                          handler: GertToolHandler)
    
    // Sync completed runs to server
    public static func syncRuns() async throws
}
\end{minted}

\subsection{Android SDK API Surface}

\begin{minted}{java}
public class GertSDK {
    // Initialise SDK with configuration
    public static void configure(GertConfig config);
    
    // Load a kit bundle from filesystem or network
    public static void loadKit(String bundlePath) 
                               throws Exception;
    
    // List all runbooks from loaded kits
    public static List<RunbookCatalogEntry> listRunbooks();
    
    // Execute a runbook
    public static GertRunResult executeRunbook(
        String id, 
        Map<String, Object> inputs,
        GertExecutionDelegate delegate
    ) throws Exception;
    
    // Resume a suspended run
    public static GertRunResult resumeRun(String runId) 
                                throws Exception;
    
    // Register a native tool handler
    public static void registerToolHandler(String name, 
                                          GertToolHandler handler);
    
    // Sync completed runs to server
    public static void syncRuns() throws Exception;
}
\end{minted}

\subsection{Execution Delegate}

The \texttt{GertExecutionDelegate} interface allows the application layer to observe run progress, handle manual prompts, and update UI.

\begin{minted}{go}
public protocol GertExecutionDelegate {
    func onStepStarted(stepId: String, stepName: String)
    func onStepCompleted(stepId: String, captures: [String: Any])
    func onStepFailed(stepId: String, error: Error)
    func onManualPrompt(prompts: [Prompt]) 
                       async -> [String: EvidenceValue]
    func onRunCompleted(runId: String, result: GertRunResult)
}
\end{minted}

The \texttt{onManualPrompt} callback is invoked when a \texttt{manual} step requires operator input. The delegate is responsible for presenting the prompt UI and collecting responses.

% ------------------------------------------------------------------ %
\section{New Capability Tokens}
% ------------------------------------------------------------------ %

The mobile execution model introduces the following capability tokens (see Section~\ref{sec:tool-capability-gate}):

\begin{center}
\begin{tabular}{ll}
\hline
\textbf{Capability} & \textbf{Description} \\
\hline
\texttt{capability/camera} & Access to device camera for photo evidence \\
\texttt{capability/location} & Access to GPS location services \\
\texttt{capability/nfc} & Access to NFC reader (equipment tag scanning) \\
\texttt{capability/biometrics} & Access to Touch ID / Face ID / fingerprint \\
\texttt{capability/bluetooth} & Access to Bluetooth LE scanning \\
\texttt{capability/notifications} & Send local push notifications \\
\hline
\end{tabular}
\end{center}

Tools that require these capabilities must declare them in the \texttt{governance.requires-capabilities} field. The SDK enforces capability grants at the application level (iOS entitlements, Android permissions).

% ------------------------------------------------------------------ %
\section{What Stays Unchanged}
% ------------------------------------------------------------------ %

The mobile execution model is explicitly designed to preserve compatibility with the core gert v2 specification. The following components are \textbf{unchanged}:

\begin{itemize}
  \item \textbf{Runbook YAML format} — \texttt{apiVersion: runbook/v2} documents are identical across CLI, server, and mobile platforms. The same runbook source file compiles to all targets.
  \item \textbf{Trace JSONL format} — Mobile runs produce the same \texttt{trace.jsonl} format described in Chapter~\ref{ch:evidence-tracing-resumption}. Event kinds, envelopes, and payload schemas are identical.
  \item \textbf{Evidence and attachment model} — SHA256 content-addressed attachments, evidence capture, and redaction rules are unchanged.
  \item \textbf{Run directory structure} — The \texttt{.runbook/runs/<run-id>/} directory layout is identical. Mobile runs can be copied to a desktop environment and resumed via \texttt{gert exec --resume}.
  \item \textbf{State snapshot contract} — The \texttt{RunState} JSON schema and checkpoint semantics are unchanged. Mobile runs are resumable on desktop and vice versa (subject to tool availability).
  \item \textbf{Governance model} — Capability gates, redaction rules, approval gates, and policy enforcement are identical across platforms.
\end{itemize}

The \texttt{client} field in the \texttt{run/started} event (see Section~\ref{sec:run-started-event-client-field}) identifies which executor produced the run. This is the only trace-level difference between CLI, server, and mobile executions.

% ------------------------------------------------------------------ %
\section{Implementation Notes}
% ------------------------------------------------------------------ %

\subsection{Cross-Platform Core}

The gert engine core (parser, planner, state machine, governance enforcer) is written in Go and compiled to native libraries using \texttt{gomobile}:

\begin{verbatim}
gomobile bind -target ios     -o GertSDK.xcframework  ./sdk/mobile
gomobile bind -target android -o gertsdk.aar          ./sdk/mobile
\end{verbatim}

The \texttt{sdk/mobile} package exposes a stable C-compatible API boundary that the Swift and Kotlin wrappers call into. This ensures that engine behaviour is identical across all platforms.

\subsection{Tool Dispatch Adapter}

The mobile SDK includes a tool dispatch adapter that intercepts tool invocations and routes them to registered native handlers. The adapter conforms to the same \texttt{ToolInvoker} interface used by the CLI executor (see Section~\ref{sec:tool-transport}).

When a runbook step calls a tool:

\begin{enumerate}
  \item The engine checks the tool's \texttt{transport} field.
  \item If \texttt{transport: native-sdk}, the adapter looks up the handler in the registry.
  \item The adapter marshals \texttt{args} and \texttt{context} to the platform's native types (Swift dictionary, Kotlin map).
  \item The handler executes and returns a \texttt{GertToolResult}.
  \item The adapter unmarshals the result back to the engine's internal representation.
\end{enumerate}

This boundary allows platform-specific tool implementations (camera, GPS, NFC) to integrate seamlessly with the core engine.

\subsection{Security Model}

Kit bundles are signed using Ed25519 signatures. The SDK verifies signatures before loading kits:

\begin{enumerate}
  \item Compute SHA256 of all files in the bundle (excluding \texttt{manifest.json}).
  \item Extract the public key identified by \texttt{signature.keyId} from the SDK's trusted key store.
  \item Verify the signature in \texttt{manifest.json} against the computed checksum.
  \item Reject the kit if verification fails.
\end{enumerate}

The application layer configures the trusted key store at SDK initialisation. Organisations deploying internal kits provide their signing keys via the SDK configuration.
